%% kapa: exam=1 round=34 year=2023 period=2 form=A subject=accounting subject_ko=회계학 seq=15 number=55
%% source: Q-Net 감정평가사 기출문제 (cst00309) artlSeq=5211789 file=raw/1008/5211789/2023년도 제34회 감정평가사 1차시험 2교시.pdf
%% conversion: kordoc 4.24.51 -> tools/gen1.py (auto); review=proofread-sonnet; flags=none
\begin{question}{55}
20×1년부터 (주)감평은 제품판매 후 2년 동안 제품하자보증을 실시하고 있다. 20×2년도에 판매된 제품에 대하여 경미한 결함은 ￦100, 치명적인 결함은 ￦4,000의 수리비용이 발생한다. 과거 경험에 따르면 10\%는 경미한 결함이, 5\%는 치명적인 결함이 발생할 것으로 예상된다. 20×1년 말에 제품보증충당부채 잔액은 ￦200이다. 20×2년 기중에 20×1년 판매된 제품에 대한 수리비용이 ￦300 지출되었다면, (주)감평의 20×2년도 재무제표에 보고할 제품보증비와 제품보증충당부채는?
\par{\small\centering\begin{tabular}{@{}cll@{\hspace{3em}}cll@{}}
 & \uline{제품보증비} & \uline{제품보증충당부채} & & \uline{제품보증비} & \uline{제품보증충당부채} \\[2pt]
① & ￦100 & ￦310 & ② & ￦210 & ￦210 \\[2pt]
③ & ￦210 & ￦310 & ④ & ￦310 & ￦210 \\[2pt]
⑤ & ￦310 & ￦310 & & & \\
\end{tabular}\par}
\end{question}
\answer{4}
